\documentclass[12pt,a4paper]{article}
\usepackage[utf8]{inputenc}
\usepackage[T1]{fontenc}
\usepackage[brazilian]{babel}
\usepackage[left=3cm,top=3cm,right=2cm,bottom=2cm]{geometry}
\usepackage{amsmath, amssymb}
\usepackage{setspace}
\usepackage[colorlinks=true, urlcolor=blue, linkcolor=blue]{hyperref}
\usepackage{xurl}
\usepackage{booktabs}
\usepackage{float}
\usepackage{minted}
\usepackage{dirtree}

\setminted{linenos, frame=single, fontsize=\small, breaklines}
\onehalfspacing

\begin{document}

\begin{titlepage}
    \begin{center}
        {\large \textbf{UNIVERSIDADE DE SÃO PAULO}} \\
        {\large \textbf{INSTITUTO DE MATEMÁTICA, ESTATÍSTICA E CIÊNCIA DA COMPUTAÇÃO}} \\

        \vspace{5cm}
        \vfill

        \begin{minipage}{0.8\textwidth}
            \centering
            {\Large \textbf{MAC0239 - EP1 - Relatório}} \\
            \vspace{0.8cm}
            {\large
            Felipe Brancalhão Curcio (16985106),\\
            Gabriela Kwong Yamaguti (16967502),\\
            Isabela Suemi Aoki (15637975),\\
            Jaime Cavalcanti Guedes (16875204),\\
            Murilo Santos de Castro (17015895),\\
            Yoonsik Hur (16905558)}
        \end{minipage}

        \vfill
        \vfill

        {\large São Paulo} \\
        {\large 2026}
    \end{center}

    \vspace*{\fill}
\end{titlepage}

\newpage
\tableofcontents
\thispagestyle{empty}

\newpage
\pagenumbering{arabic}

% 1. Introdução -----------------------------------------------------
\section{Introdução}

O presente trabalho teve como objetivo a criação de um provador de sequentes $A_1, \dots, A_n \vdash B$ do cálculo proposicional por \textit{tableaux} semântico, mostrando as regras \textit{alpha} e \textit{beta} aplicadas. Para cada sequente, o programa indica se ele é válido e, caso não seja, mostra um contra-exemplo.

\noindent O método usado é o de refutação: as premissas são marcadas inicialmente como verdadeiras ($T$) e a conclusão como falsa ($F$), de forma que o programa tenta construir uma valoração que torne as premissas verdadeiras e a conclusão falsa. Se todos os ramos fecham, essa valoração não existe e, portanto, o sequente é válido. Caso algum ramo termine aberto e totalmente expandido, ele fornece um contra-exemplo.

\noindent A linguagem aceita contém os conectivos $\neg$, $\wedge$, $\vee$ e $\rightarrow$, escritos em Python como \texttt{\textasciitilde}, \texttt{\&}, \texttt{|} e \texttt{>>}.

% 2. Desenvolvimento ------------------------------------------------
\section{Desenvolvimento}

% 2.1. Como utilizar o script ---------------------------------------
\subsection{Como utilizar o script}

Para usar o provador, é necessário ter Python 3 instalado (testado com Python 3.13). Não há outras dependências externas. Além disso, usamos a seguinte estrutura de diretórios (com uma estrutura diferente, as importações podem ficar incorretas):

\dirtree{%
.1 meu\_projeto.
.2 tableaux.py \hfill \texttt{\# implementação do provador}.
.2 testes.py \hfill \texttt{\# implementação dos testes}.
.2 output.txt \hfill \texttt{\# saída dos testes}.
.2 arquivos\_enunciado.
.3 code\_base\_EP1.py \hfill \texttt{\# código-base fornecido}.
.3 code\_base\_EP1\_demo.py \hfill \texttt{\# exemplos de uso do código-base}.
.2 relatorio.
.3 relatorio.pdf \hfill \texttt{\ # pdf do relatório}.
}

\bigskip
\noindent Para testar o provador, há duas maneiras:

\begin{enumerate}
    \item Criar um novo script que chame as funções disponíveis em \texttt{tableaux.py}, como no exemplo abaixo:

\begin{minted}{python}
from tableaux import *

# todos os átomos usados precisam ser criados
P, Q, R = vars('P', 'Q', 'R')

imprimir_resultado(ArgumentForm(P, P >> Q, conclusion=Q))   # P, P -> Q |- Q
imprimir_resultado(ArgumentForm(P >> Q, Q, conclusion=P))   # P -> Q, Q |- P
imprimir_resultado(ArgumentForm(conclusion=P | ~P))         # |- P v ~P (sem premissas)
\end{minted}

    Saída dos dois primeiros sequentes (os símbolos $\vdash$, $\beta$ e $\rightarrow$ são os que o programa imprime):

\begin{minted}[escapeinside=||]{text}
Sequente: P, P => Q |$\vdash$| Q
Resultado: VÁLIDO
Regras: [('|$\beta\rightarrow$|', (T,P => Q))]

Sequente: P => Q, Q |$\vdash$| P
Resultado: INVÁLIDO
Contra-exemplo: P = False, Q = True
Regras: [('|$\beta\rightarrow$|', (T,P => Q))]
\end{minted}

    \item Rodar os testes unitários com \texttt{python tableaux.py}.
\end{enumerate}

\paragraph{Conectivos.}
\begin{table}[H]
\centering
\begin{tabular}{@{}ll@{}}
\toprule
\textbf{Notação da aula} & \textbf{No Python} \\
\midrule
$\neg A$ & \texttt{\textasciitilde A} \\
$A \wedge B$ & \texttt{A \& B} \\
$A \vee B$ & \texttt{A | B} \\
$A \rightarrow B$ & \texttt{A >> B} \\
\bottomrule
\end{tabular}
\end{table}

\paragraph{Cuidados ao escrever sequentes:}
\begin{itemize}
    \item No Python, \texttt{>>} tem prioridade sobre \texttt{\&} e \texttt{|}. Assim, \texttt{P \& Q >> R} significa \texttt{P \& (Q >> R)}, e não \texttt{(P \& Q) >> R}. Deve-se colocar parênteses em volta de toda implicação.
    \item \texttt{R = vars('R')} devolve um gerador, e não um átomo. Deve-se usar \texttt{R, = vars('R')} (com a vírgula) ou criar todos os átomos de uma vez.
    \item Apenas \texttt{\textasciitilde}, \texttt{\&}, \texttt{|} e \texttt{>>} fazem parte da linguagem do enunciado. Outros tipos do código-base (\texttt{<=>}, \texttt{<<} e constantes) geram \texttt{TypeError}.
\end{itemize}

% 2.2. Estrutura interna do programa ---------------------------------
\subsection{Estrutura interna do programa}

O programa \texttt{tableaux.py} disponibiliza três funções ao usuário:
\begin{itemize}
    \item \texttt{imprimir\_resultado(sequente)}: imprime o sequente, se ele é válido, o contra-exemplo (se houver) e as regras usadas;

    \item \texttt{provar(sequente)}: retorna a tripla \texttt{(valido, ramo, regras)}, em que \texttt{valido} é do tipo booleano (\texttt{True}/\texttt{False}), \texttt{ramo} é o ramo aberto (ou \texttt{None}, se o sequente é válido) e \texttt{regras} é a lista \texttt{[(regra, fórmula marcada), ...]}. O funcionamento segue os passos:
    \begin{itemize}
        \item Cria o ramo inicial, com as premissas marcadas como verdadeiras e a conclusão marcada como falsa, e o coloca na pilha de ramos;
        \item Em cada iteração, é retirado um ramo da pilha e, se ele estiver fechado, é descartado;
        \item Chama \texttt{escolher()} para obter a próxima fórmula a expandir. Se não houver nenhuma, o ramo está aberto e totalmente expandido, então o sequente é inválido e esse ramo é devolvido como contra-exemplo;
        \item Caso contrário, classifica a fórmula, registra a regra em \texttt{regras} e aplica a regra: se for $\alpha$, adiciona os componentes ao próprio ramo e o devolve à pilha; se for $\beta$, cria uma cópia do ramo para cada componente, adiciona o componente à cópia e empilha as cópias;
        \item Se a pilha esvazia, todos os ramos fecharam e o sequente é válido.
    \end{itemize}

    \item \texttt{contra\_exemplo(sequente, ramo)}: retorna um dicionário \texttt{\{átomo: valor\}} com todos os átomos do sequente, somente para sequentes inválidos. Átomos que não aparecem no ramo aberto recebem \texttt{True} (valor livre).
\end{itemize}

\noindent Internamente, a função \texttt{classificar\_tipo\_regra\_comp(FormulaMarcada)} recebe uma fórmula marcada e retorna a tripla \texttt{(tipo, regra, componentes)}, em que:
\begin{itemize}
    \item \texttt{tipo} indica como usar os componentes: \texttt{alpha} (todos valem ao mesmo tempo, no mesmo ramo), \texttt{beta} (o ramo se divide, um para cada componente) ou \texttt{literal} (átomo, que não pode ser decomposto);
    \item \texttt{regra} é o rótulo que aparece na saída, como \texttt{'$\alpha\rightarrow$'};
    \item \texttt{componentes} é a lista de novas fórmulas marcadas produzidas pela regra.
\end{itemize}

\noindent Além disso, o programa conta com as seguintes estruturas:
\begin{itemize}
    \item \texttt{FormulaMarcada}: par \texttt{(sinal, fórmula)}, em que \texttt{True} representa $T$ e \texttt{False} representa $F$. Inicialmente, cada premissa é $(T, \text{premissa})$ e a conclusão é $(F, \text{conclusão})$.

    \item \texttt{Ramo}: guarda a lista de fórmulas marcadas compostas ainda pendentes (\texttt{pendentes}), o dicionário \texttt{\{átomo: sinal\}} com os literais encontrados (\texttt{literais}) e o booleano \texttt{fechado} (\texttt{True} se o ramo contém uma contradição). Essa classe possui três métodos principais:
    \begin{itemize}
        \item \texttt{adicionar(FormulaMarcada)}: separa o que entra. Se for um átomo, vai para \texttt{literais} e, se esse átomo já existir com o sinal oposto, o ramo é fechado. Se for composta, vai para \texttt{pendentes};
        \item \texttt{copiar()}: cria um novo ramo com lista e dicionário próprios, de modo que alterar a cópia não afeta o original;
        \item \texttt{escolher()}: remove e devolve o próximo pendente a expandir, preferindo as regras $\alpha$. Se não há pendentes, devolve \texttt{None}.
    \end{itemize}

    \item Pilha de ramos: lista usada como pilha (LIFO) em \texttt{provar}, que guarda os ramos ainda não examinados e resulta em uma busca em profundidade.

    \item Lista de regras aplicadas: guarda os pares \texttt{(regra, fórmula marcada)} na ordem de aplicação, cobrindo todos os ramos visitados.
\end{itemize}

% 2.3. Sequentes para testar o programa ------------------------------
\subsection{Sequentes para testar o programa}
Para testar o provador, usamos o script \texttt{testes.py}, com 20 sequentes: 14 válidos e 6 inválidos. Para rodar os testes, basta executar \texttt{python3 testes.py} no diretório do projeto. Os sequentes usados foram divididos nos seguintes grupos:

\paragraph{Argumentos clássicos e conjunção com três filhos (válidos):}
\begin{enumerate}
    \item Modus Ponens: $P,\ P \rightarrow Q \vdash Q$.
    \item Modus Tollens: $P \rightarrow Q,\ \neg Q \vdash \neg P$.
    \item Silogismo Disjuntivo: $P \vee Q,\ \neg P \vdash Q$.
    \item Silogismo Hipotético: $P \rightarrow Q,\ Q \rightarrow R \vdash P \rightarrow R$.
    \item Lei de De Morgan: $\neg(P \wedge Q) \vdash \neg P \vee \neg Q$.
    \item Conjunção com três filhos: $P \wedge Q \wedge R \vdash Q \wedge R$.
\end{enumerate}

\paragraph{Teoremas sem premissas (válidos):}
\begin{enumerate}
    \item Terceiro Excluído: $\vdash P \vee \neg P$.
    \item Lei de Peirce: $\vdash ((P \rightarrow Q) \rightarrow P) \rightarrow P$.
    \item Distributividade da disjunção sobre a conjunção:
    $\vdash (P \vee (Q \wedge R)) \rightarrow ((P \vee Q) \wedge (P \vee R))$.
\end{enumerate}

\paragraph{Casos especiais e fórmulas compostas (válidos):}
\begin{enumerate}
    \item Identidade: $P \vdash P$, com fechamento imediato, sem aplicação de regras de expansão.
    \item Princípio da Explosão: $P,\ \neg P \vdash Q$.
    \item Dupla Negação: $\neg\neg P \vdash P$.
    \item Disjunção com três filhos e aplicação de $\alpha$ antes de $\beta$:
    $P \vee Q \vee R,\ (\neg P \wedge \neg Q) \vdash R$.
    \item Fórmula composta com quatro átomos:
    $((A \vee B) \rightarrow Q) \wedge \neg S,\ A \vdash Q \wedge \neg S$.
\end{enumerate}

\paragraph{Falácias e sequentes inválidos:}
\begin{enumerate}
    \item Afirmação do Consequente: $P \rightarrow Q,\ Q \vdash P$.
    \item Negação do Antecedente: $P \rightarrow Q,\ \neg P \vdash \neg Q$.
    \item Non Sequitur: $P \vdash Q$, com ramo aberto sem aplicação de regras de expansão.
    \item Inferência indevida a partir de uma disjunção: $P \vee Q,\ P \vdash \neg Q$.
    \item Contra-exemplo com átomo livre: $P \vee Q \vdash P \wedge R$.
    \item Fórmula composta com quatro átomos:
    $((A \vee B) \rightarrow Q) \wedge \neg S \vdash Q \wedge \neg S$.
\end{enumerate}

\paragraph{Verificação dos resultados:} para cada sequente, o script compara o resultado do provador com o resultado esperado e com a tabela-verdade do código-base. Para os sequentes inválidos, também verifica se o contra-exemplo torna todas as premissas verdadeiras e a conclusão falsa.

\paragraph{Fórmulas fora da linguagem:} também testamos o bicondicional \texttt{P.iff(Q)}, a implicação reversa \texttt{P << Q} e a constante \texttt{True}. Nos três casos, o script verifica se o provador gera \texttt{TypeError}, pois essas fórmulas não fazem parte da linguagem do enunciado.

\paragraph{Resultados obtidos:} os 20 sequentes passaram nos testes, incluindo a verificação dos contra-exemplos. Os três casos fora da linguagem geraram \texttt{TypeError} e todas as oito regras de expansão foram usadas.

% 2.4. Estratégias usadas na implementação ---------------------------
\subsection{Estratégias usadas na implementação}

\paragraph{Refutação:} para decidir $A_1,\dots,A_n \vdash B$, o programa marca as premissas como $T$ e a conclusão como $F$, procurando uma valoração para essas fórmulas. Se todos os ramos fecham, o sequente é válido. Caso algum ramo termine aberto e totalmente expandido, o sequente é inválido, gerando o contra-exemplo.

\paragraph{Escolha da regra em um ramo:} o método \texttt{escolher} aplica as regras $\alpha$ antes das $\beta$. Como as regras $\alpha$ não ramificam, adiar as regras $\beta$ evita expandir as mesmas fórmulas em vários ramos. Assim, muitos ramos fecham antes de se dividirem.

\paragraph{Escolha do ramo:} os ramos pendentes ficam em uma pilha (LIFO), o que equivale a uma busca em profundidade. O algoritmo para no primeiro ramo aberto e completamente expandido, pois já basta ele para exibir um contra-exemplo.

\paragraph{Fechamento do ramo:} ao adicionar um literal, verifica-se se o mesmo átomo já está no ramo com o sinal oposto. Os literais ficam em um dicionário \texttt{\{átomo: sinal\}}, o que torna essa verificação imediata. Fórmulas compostas contraditórias fecham o ramo apenas depois de expandidas.

\paragraph{Negação:} na tabela do enunciado, $F\ \neg C$ tem $\beta_1=\beta_2=T\ C$. Dividir o ramo criaria dois ramos idênticos, então a negação é tratada como $\alpha$ nos dois sinais. O nome \texttt{$\beta\neg$} é mantido no registro de regras para coincidir com a tabela.

\paragraph{Conjunções e disjunções com $n$ filhos:} o código-base transforma \texttt{P \& Q \& R} em um único \texttt{And} com três filhos. As regras foram generalizadas: $T\ (A_1\wedge\dots\wedge A_n)$ adiciona os $n$ componentes ao mesmo ramo, e $F\ (A_1\wedge\dots\wedge A_n)$ abre $n$ ramos. O mesmo vale para $\vee$, com os papéis de $\alpha$ e $\beta$ trocados.

\paragraph{Ramificação:} em uma regra $\beta$, cada componente é adicionado a uma \textbf{cópia} do ramo atual, com lista e dicionário próprios. Assim, alterar um ramo não afeta os outros.

\paragraph{Contra-exemplo:} átomos do sequente que não aparecem no ramo não influenciam a refutação, então recebem \texttt{True} e são exibidos como \textit{valor livre}.

\paragraph{Registro de regras:} uma lista única guarda pares \texttt{(regra, fórmula marcada)} na ordem de aplicação, cobrindo todos os ramos visitados. Em sequentes inválidos, ela termina no primeiro ramo aberto encontrado.

% 2.5. Dificuldades encontradas --------------------------------------
\subsection{Dificuldades encontradas}
\begin{itemize}
    \item \textbf{Precedência de operadores no Python:} o operador \texttt{>>} tem prioridade maior que \texttt{\&} e \texttt{|}, então \texttt{P \& Q >> R} significa \texttt{P \& (Q >> R)}, e não \texttt{(P \& Q) >> R}. Isso gerou sequentes diferentes dos pretendidos. Resolvemos exigindo parênteses em toda implicação e criando um teste para esse caso.
    \item \textbf{Criação de um único átomo:} \texttt{R = vars('R')} devolve um gerador, não um átomo. É preciso escrever \texttt{R, = vars('R')} ou criar todos os átomos de uma vez.
    \item \textbf{Regra $\beta\neg$:} decidir como representar $F\ \neg C$. Dividir em dois ramos idênticos seria redundante, então o tipo foi tratado como $\alpha$ e o nome da regra foi mantido.
\end{itemize}

% 3. Conclusão --------------------------------------------------------
\section{Conclusão}
O provador implementa o método de \textit{tableaux} semântico para lógica proposicional, decidindo a validade de qualquer sequente da linguagem aceita, exibindo as regras aplicadas e, quando o sequente é inválido, apresentando um contra-exemplo. A priorização das regras $\alpha$ e a busca em profundidade ajudam na eficiência do método.

\end{document}